\documentclass[12pt]{article}
\usepackage[T1]{fontenc}
\usepackage[utf8]{inputenc}
\usepackage{lmodern}
\usepackage{textcomp}
\usepackage{enumitem}
\usepackage{geometry}
\usepackage{graphicx}
\usepackage{amsmath, amssymb}
\geometry{margin=1in}
\begin{document}
\begin{center}
Chemistry - Graduate Level Assessment 20 \
Total Marks: 40 \
Answer all questions.
\end{center}

\section*{Multiple Choice (10 marks)}
\begin{enumerate}[label=\arabic*.]
\item What is the maximum number of phases that can be at equilibrium with each other in a three component mixture?
\begin{enumerate}[label=(\alph*)]
    \item 6
    \item 9
    \item 2
    \item 4
    \item 10
\end{enumerate}
\item Calculate the de Broglie wavelength for (a) an electron with a kinetic energy of $100 \mathrm{eV}$
\begin{enumerate}[label=(\alph*)]
    \item 0.175 nm
    \item 0.080 nm
    \item 0.098 nm
    \item 0.123 nm
    \item 0.135 nm
\end{enumerate}
\item One cubic millimeter of oil is spread on the surface of water so that the oil film has an area of 1 square meter. What is the thickness of the oil film in angstrom units?
\begin{enumerate}[label=(\alph*)]
    \item 200 \ensuremath{\AA}
    \item 1 \ensuremath{\AA}
    \item 0.1 \ensuremath{\AA}
    \item 50 \ensuremath{\AA}
    \item 500 \ensuremath{\AA}
\end{enumerate}
\item Which of the following molecules is a strong electrolyte when dissolved in water?
\begin{enumerate}[label=(\alph*)]
    \item CH 3 COOH
    \item CO 2
    \item O$_{2}$
    \item H$_{2}$O
    \item NH 3
\end{enumerate}
\item Which of the following is an n-type semiconductor?
\begin{enumerate}[label=(\alph*)]
    \item Silicon
    \item Boron-doped silicon
    \item Pure germanium
    \item Arsenic-doped silicon
    \item Indium Phosphide
\end{enumerate}
\end{enumerate}

\section*{True / False (10 marks)}
\textit{Answer True or False}
\begin{enumerate}[label=\arabic*.]
\setcounter{enumi}{5}
\item True or False: The nuclide 31 P has an NMR frequency of 115.5 MHz in a 20.0 T magnetic field.
\item True or False: The temperature change of a diatomic gas can be directly calculated as 13.5 K from the given values of q and w.
\item True or False: When two coins are tossed, the expected number of heads, denoted as ⟨w⟩, is 1.
\item True or False: Element X has a valence of +4, indicating it can accept four electrons in reactions.
\item True or False: The carbon configuration $1 s^2 2 s^2 2 p^2$ corresponds to a total of 20 microstates in quantum mechanics.
\end{enumerate}

\section*{Long Form Response (20 marks)}
\textit{Answer each question with a clear and complete explanation.}
\begin{enumerate}[label=\arabic*.]
\setcounter{enumi}{10}
\item Discuss the kinetic energy of a particle with a mass of 1.0 g released from rest near the Earth's surface, where g = 8.91 m/$s^2$, after 1.0 s. Provide a detailed analysis and calculations.
\item Discuss the total binding energy of the carbon-12 nucleus and analyze the significance of the average binding energy per nucleon in the context of nuclear stability.
\end{enumerate}

\vfill
\noindent\textit{End of Paper}
\end{document}